\documentclass[11pt,a4paper]{article}
\usepackage[T1]{fontenc}
\usepackage[utf8]{inputenc}
\usepackage{lmodern}
\usepackage[margin=2.2cm]{geometry}
\usepackage{xcolor}
\usepackage{listings}
\usepackage{enumitem}
\usepackage{hyperref}
\hypersetup{colorlinks=true,urlcolor=blue!50!black,linkcolor=black}
\lstset{basicstyle=\ttfamily\small,backgroundcolor=\color{gray!10},frame=single,rulecolor=\color{gray!40},
        columns=fullflexible,breaklines=true,xleftmargin=4pt,xrightmargin=4pt}
\setlist{itemsep=2pt,topsep=4pt}
\newcommand{\f}[1]{\texttt{#1}}
\title{Room estimation with two rooms on the real robot\\[2pt]\large How to run it and how to make the robot cross doors}
\author{For Noé, prepared with Claude for Pablo}
\date{26 November 2026}
\sloppy
\begin{document}
\maketitle

\noindent All paths are relative to the repository root \f{robocomp/\allowbreak{}components/\allowbreak{}active\_\allowbreak{}inference}
(branch \f{main}). This note is meant to be read by you and by your Claude session.

\section{The idea in one paragraph}
The robot starts in room~1 with an empty memory, estimates the room's layout while driving, finds the door, crosses
it, estimates room~2, and crosses back. The cortex fleet that does this is the same one used in Webots:
\f{retina}, \f{room\_\allowbreak{}concept}, \f{residual\_\allowbreak{}concept}, \f{controller}, \f{door\_\allowbreak{}concept}, \f{ltsm\_\allowbreak{}agent}
(started by \f{room\_\allowbreak{}concept/\allowbreak{}estimation\_\allowbreak{}scripts/\allowbreak{}scripts/\allowbreak{}run\_\allowbreak{}crossing.sh}). Only the simulator-specific steps of the
scripts have to be removed.

\section{How the robot is told to cross a door}
The controller watches a file, \f{controller/\allowbreak{}remote\_\allowbreak{}cmd.txt}. It reads it every 300\,ms and deletes it, so each command
runs exactly once (code: \f{controller/\allowbreak{}src/\allowbreak{}specificworker.cpp}, the ``REMOTE RUN/STOP'' timer). There are two ways to
use it.

\subsection*{Option 1 (recommended): the robot runs the tour itself}
\begin{lstlisting}
echo "explore_tour on" > controller/remote_cmd.txt
\end{lstlisting}
The exploration drive (\f{controller/\allowbreak{}src/\allowbreak{}exploration\_\allowbreak{}drive.h}, spec
\f{docs/\allowbreak{}superpowers/\allowbreak{}specs/\allowbreak{}2026-11-05-exploration-tour-design.md}) explores the current room until the room is
confirmed, crosses the nearest door it has not crossed yet, explores the next room, and, when no uncrossed door is left,
returns through the door it came in by. With two rooms this is exactly room~1 $\to$ cross $\to$ room~2 $\to$ cross back.
It prints \f{[explore-tour] DONE} in the controller log when it finishes. The reference script is
\f{room\_\allowbreak{}concept/\allowbreak{}estimation\_\allowbreak{}scripts/\allowbreak{}scripts/\allowbreak{}explore\_\allowbreak{}tour.sh} (Estimate mode, empty memory, recognition on).
In the Webots apartment this passed 6 of 6 tours.

\subsection*{Option 2: scripted legs}
Write a mission script into the same file:
\begin{lstlisting}
mission "leg 1"
  open new
  cross new retry=2
\end{lstlisting}
\f{cross new} means ``the door whose far side I do not remember yet'', so no door name is needed in advance. After the
first crossing the door has a name, and \f{cross room\_\allowbreak{}1} brings the robot back. A single line \f{cross door\_\allowbreak{}N} also
works when you know the door node's name. \f{warm\_\allowbreak{}start\_\allowbreak{}tour.sh} shows how legs are sent and how their end is
detected from \f{controller.log} (\f{passage DONE}, \f{mission 'leg N': completed /\allowbreak{} failed}).

\subsection*{Other useful commands (same file)}
\f{run}, \f{stop}, \f{cancel}, \f{pause}, \f{unpause}, \f{move <m>}, \f{turn <deg>}, \f{explore\_\allowbreak{}tour off}.
\f{stop} is the remote equivalent of the Stop button in the controller window.

\section{What is different on the real robot}
\begin{enumerate}
\item \textbf{Open the doors by hand before the run.} The \f{open} step goes through the DoorControl interface
  (\f{door\_\allowbreak{}concept/\allowbreak{}src/\allowbreak{}door\_\allowbreak{}actuator.h}). It is provider-agnostic, but today the only provider is the Webots
  supervisor, and its door matching uses simulator-only ground-truth poses. On the robot nobody serves it. An
  \f{open} step is satisfied when the door is \emph{seen} open, so prop the door fully open and leave it open.
\item \textbf{Remove the Webots-only steps from the scripts:}
  \begin{itemize}
  \item \f{tp.py} (teleports the robot to the start pose): place the robot by hand instead and use
    \f{run\_\allowbreak{}crossing.sh <dir> here}, which does not teleport;
  \item \f{door.py} (opens/closes the simulated doors);
  \item the ground-truth pose log \f{gt\_\allowbreak{}pose.tsv}, \f{true\_\allowbreak{}room.py}, \f{doors\_\allowbreak{}crossed.py}: they read the simulator's pose and cannot work on hardware. (\f{here} still calls
    \f{tp.py get} once; on the robot that line only logs an error.) \f{loc\_\allowbreak{}ok.py} is fine: it reads
    \f{room\_\allowbreak{}concept}'s own scan residual, not ground truth, and is a useful ``is the pose good enough to start'' check.
  \end{itemize}
  Which doors were crossed is then judged from the robot's own logs plus a person watching.
\item \textbf{Ground truth for the rooms.} As agreed in \f{REAL\_\allowbreak{}ROBOT\_\allowbreak{}VALIDATION\_\allowbreak{}PLAN.md}
  (in Pablo's \f{Papers Ongoing/\allowbreak{}2026/\allowbreak{}ICRA-rejected} folder): survey each room with a laser meter. The meter gives
  distances only, so the corners are computed by least squares. Not built yet: the survey file format, the
  least-squares fit, and a scorer that compares the learnt layout (\f{layouts/\allowbreak{}learnt/\allowbreak{}*.svg}, \f{ltsm\_\allowbreak{}history}) with it.
\item \textbf{Safety.} Stop with \f{echo stop > controller/\allowbreak{}remote\_\allowbreak{}cmd.txt} or the Stop button. Keep someone with a
  joystick or emergency stop next to the robot. Stop agents with SIGTERM (\f{kill}), never \f{kill -9}: the agents
  clean their nodes from the shared graph only on a graceful exit.
\end{enumerate}

\section{Suggested steps}
\begin{enumerate}
\item Write \f{real\_\allowbreak{}two\_\allowbreak{}rooms.sh}: a copy of \f{explore\_\allowbreak{}tour.sh} without the simulator steps listed above. Keep the
  parts that matter: archive (never delete) the old memory (\f{ltsm\_\allowbreak{}agent/\allowbreak{}etc/\allowbreak{}ltsm\_\allowbreak{}memory.json},
  \f{door\_\allowbreak{}concept/\allowbreak{}etc/\allowbreak{}door\_\allowbreak{}identities.csv}, \f{ltsm\_\allowbreak{}history}), set \f{MapMode = "estimate"} in
  \f{room\_\allowbreak{}concept/\allowbreak{}etc/\allowbreak{}config.toml} and \f{Recognition.enabled = true} in \f{ltsm\_\allowbreak{}agent/\allowbreak{}etc/\allowbreak{}config.toml} (restored on
  exit), start the fleet with \f{run\_\allowbreak{}crossing.sh <dir> here}, send \f{explore\_\allowbreak{}tour on}, wait for
  \f{[explore-tour] DONE}, then \f{run\_\allowbreak{}crossing.sh <dir> collect} to copy the logs and stop the fleet.
\item Dry-run that exact script in Webots on the two-room world (\f{webots-shadow/\allowbreak{}worlds/\allowbreak{}piso/\allowbreak{}tworooms-piso-door1open.wbt}) with the door left
  open, so the only thing that changes on the robot is the hardware.
\item On the robot: first a short run with one room and one crossing (option~2, a single \f{cross new} leg), then the
  full tour (option~1).
\item Survey both rooms with the laser meter and score the learnt layouts.
\end{enumerate}

\section{Known unknowns}
\begin{itemize}
\item Nobody has checked how the exploration drive copes with real sensor noise, or with a real doorway narrower than
  the simulated one. The Webots dry run will not show this; the first robot run is the real test.
\item The controller runs the three-mode hybrid (merged in main 37444cd7, \f{[Policy] HybridChain = true} with the
  chain monitor on): open space is driven by the spline chain (planner, spline, tracker), tight stretches such as doorways
  are handed to the local policy, and a recovery mode re-plans or waits when progress fails. The modes switch automatically;
  nothing has to be selected. Its pose-noise inputs (\f{[Policy] SigmaRelXyM}, \f{SigmaRelYawRad}, \f{PoseDriftSigmaM},
  \f{PoseDriftTimeS}) were fitted on Webots runs, and the config says to re-fit them on the real robot
  (\f{controller/tools/pose\_error\_stats.py} needs ground truth, so this is a later step).
\item The robot starts with an empty memory (Estimate mode). The warm-start mode, which starts from a remembered map,
  accepts only \f{room\_\allowbreak{}1} as the start room for now.
\end{itemize}

\end{document}
